\chapter{数学速查：条件先于公式}
\section{整除、卷积与筛法}\index{Dirichlet convolution}\index{Mobius inversion}
对正整数定义狄利克雷卷积
\[
(f*g)(n)=\sum_{d\mid n}f(d)g(n/d).
\]
单位元为 $\varepsilon(n)=[n=1]$，常函数 $1(n)=1$，恒等函数 $\operatorname{id}(n)=n$。
关键恒等式是 $\mu*1=\varepsilon$、$\varphi*1=\operatorname{id}$。
若 $F(n)=\sum_{d\mid n}f(d)$，则 $f(n)=\sum_{d\mid n}\mu(d)F(n/d)$。
\[
\sum_{i=1}^{a}\sum_{j=1}^{b}[\gcd(i,j)=1]
=\sum_{d=1}^{\min(a,b)}\mu(d)\lfloor a/d\rfloor\lfloor b/d\rfloor.
\]
整除分块：当左端点为 $l$、商为 $q=\lfloor n/l\rfloor$，右端点为 $r=\lfloor n/q\rfloor$。
同时包含多个商时，取各自右端点的最小值。
杜教筛由卷积恒等式获得
\[
M(n)=1-\sum_{l=2}^{n}M(\lfloor n/l\rfloor),\qquad
\Phi(n)=\frac{n(n+1)}2-\sum_{l=2}^{n}\Phi(\lfloor n/l\rfloor),
\]
求和按等商区间合并；第二式求和中没有额外乘以 $l$。

\section{模运算的适用条件}\index{CRT}\index{Euler theorem}\index{Lucas theorem}
模逆元存在当且仅当 $\gcd(a,m)=1$。费马逆元 $a^{p-2}$ 仅适用于素数模数及 $a\not\equiv0$。
对正模数 $m,n$，同余式 $x\equiv a\pmod m$、$x\equiv b\pmod n$ 可合并当且仅当
$\gcd(m,n)\mid(b-a)$；新模数为 $\operatorname{lcm}(m,n)$。
\[
\binom nk\equiv\prod_i\binom{n_i}{k_i}\pmod p\quad(p\text{ 为素数，}n_i,k_i\text{ 是 }p\text{ 进制位}).
\]
Lucas 不能直接套在合数模数。阶乘与逆阶乘预处理要求上限小于模数；超过模数会出现阶乘为零。
扩展欧拉降幂中，非互质情况不能直接将指数对 $\varphi(m)$ 取模；需要判断原指数是否达到足够阈值。
大整数指数可以在读入时同时维护余数与是否超过阈值。

\section{素数模平方根}\index{Cipolla}\index{quadratic residue}
奇素数 $p$ 下，对 $n\not\equiv0\pmod p$，Euler 判别式
$n^{(p-1)/2}\equiv1\pmod p$ 当且仅当存在平方根，否则为 $-1$。
当 $p\equiv3\pmod4$ 且有解时，$r=n^{(p+1)/4}$ 为一根，另一根是 $p-r$。

Cipolla 选择 $a$ 使 $w=a^2-n$ 是非平方，令 $\omega^2=w$。
在扩域中取 $\alpha=a+\omega$，有 $\alpha^p=a-\omega$，故 $\alpha^{p+1}=n$。
令 $\beta=\alpha^{(p+1)/2}$，则 $\beta^2=n$ 且
$\beta^{p-1}=n^{(p-1)/2}=1$，所以 $\beta$ 落在原素数域中。
实现见第~\pageref{compact-mod_sqrt}~页；零和模数 2 单独处理。
模数为合数时不能照搬此判别和扩域构造。

\section{合数模数组合数与插值}\index{exLucas}\index{Lagrange interpolation}
对素数幂 $q=p^e$，记 $v_p(n!)=\sum_{j\ge1}\lfloor n/p^j\rfloor$，
$U_p(n)=n!/p^{v_p(n!)}$。设
$t=v_p(n!)-v_p(k!)-v_p((n-k)!)$。当 $t\ge e$ 时组合数模 $q$ 为零；否则
\[
\binom nk\equiv U_p(n)\,U_p(k)^{-1}\,U_p(n-k)^{-1}\,p^t\pmod q.
\]
这里仅对与 $p$ 互素的单位部分求逆。各素数幂上的结果通过 CRT 合并，
对应第~\pageref{compact-ExLucas}~页的实现。

在素数域中，$n$ 个模 $p$ 两两不同的横坐标确定唯一的次数小于 $n$ 的多项式。
定义 $w_i=y_i/\prod_{j\ne i}(x_i-x_j)$，则
\[
f(k)=\sum_{i=0}^{n-1}w_i\prod_{j<i}(k-x_j)\prod_{j>i}(k-x_j).
\]
该式不除以 $k-x_i$，查询点命中样本点时仍然成立。
当 $x_i=i$ 时，分母为 $i!\,(n-1-i)!\,(-1)^{n-1-i}$，要求 $n\le p$。
实现见第~\pageref{compact-Lagrange}~页。
对整数幂前缀和 $S_d(k)=\sum_{i=1}^k i^d$，当 $0\le d\le p-2$ 时，
它在模 $p$ 意义下可用次数至多 $d+1$ 的多项式表示；取 $0,1,\ldots,d+1$
的样本即可。不能只因某序列前几项能插值，就断言后续仍由该多项式给出。

\section{组合计数}\index{inclusion-exclusion}\index{Catalan}\index{Stirling numbers}
容斥：$|\bigcup A_i|=\sum_{\emptyset\ne S}(-1)^{|S|+1}|\bigcap_{i\in S}A_i|$。
二项式反演：若 $b_n=\sum_{k=0}^{n}\binom nk a_k$，则
$a_n=\sum_{k=0}^{n}(-1)^{n-k}\binom nk b_k$。
\[
\sum_k\binom ak\binom b{n-k}=\binom{a+b}n,\qquad
C_n=\binom{2n}n-\binom{2n}{n+1}.
\]
差分形式的 Catalan 公式避免直接除以 $n+1$，在模数下尤其有用，但组合数本身仍须正确求模。
第二类 Stirling 数满足
\[
\left\{\begin{matrix}n\\k\end{matrix}\right\}
=\left\{\begin{matrix}n-1\\k-1\end{matrix}\right\}
+k\left\{\begin{matrix}n-1\\k\end{matrix}\right\},
\quad S(0,0)=1.
\]
第一类无符号 Stirling 数把系数 $k$ 换为 $n-1$。
\begin{center}
\begin{tabular}{lll}
\toprule
球与盒 & 空盒允许 & 空盒禁止\\\midrule
不同球，不同盒 & $k^n$ & $k!S(n,k)$\\
相同球，不同盒 & $\binom{n+k-1}{k-1}$ & $\binom{n-1}{k-1}$\\
不同球，相同盒 & $\sum_{j=0}^{k}S(n,j)$ & $S(n,k)$\\
\bottomrule
\end{tabular}
\end{center}
上表默认 $k\ge1$；$n=0$、$k=0$ 等边界应按原问题约定单独处理。

\section{群作用计数}\index{Burnside}\index{Polya}
有限群 $G$ 作用在有限集合 $X$ 上时，轨道数为
\[
|X/G|=\frac1{|G|}\sum_{g\in G}|\operatorname{Fix}(g)|.
\]
对位置置换的 $q$ 色染色，$|\operatorname{Fix}(g)|=q^{c(g)}$，其中 $c(g)$ 为循环数。
环形项链只商去旋转；手链还商去反射，二者群不同，不能共用分母。
模数整除 $|G|$ 时不能使用模逆元，须先在合适的整数模数下求和再做精确除法。

\section{生成函数与递推}\index{generating functions}\index{Berlekamp-Massey}
普通生成函数 $A(x)=\sum a_nx^n$ 的乘法对应普通卷积。
指数生成函数 $A(x)=\sum a_nx^n/n!$ 的乘法对应
$c_n=\sum_k\binom nk a_kb_{n-k}$。
形式幂级数求逆要求常数项可逆；$\log A$ 取 $A(0)=1$，$\exp A$ 取 $A(0)=0$。
有限域中积分涉及 $1/i$，必须避免 $i$ 被特征整除。
\[
B_{\rm new}=B(2-AB)\pmod{x^{2m}},\qquad
(\log A)'=A'/A.
\]
BM 从样本得到最短线性递推。要对未来项使用，必须另外证明序列确实具有不超过相应阶数的线性递推；
少量数值吻合不是证明。矩阵转移、有限状态计数常可提供阶数上界。

\section{线性代数与生成树}\index{rank}\index{Matrix Tree theorem}
域上的线性方程组一致时，解集为一个特解加零空间，零空间维数为未知数个数减秩。
浮点消元中的秩依赖误差阈值；模素数消元是精确域运算；合数模数下不可随意除以非零元。
无向图 Laplacian 的对角元为度数，非对角元为负的边重数。删去同一行列后的行列式为生成树数。
自环不参与生成树；重边需要累加。只有一个顶点时空余子式行列式约定为 1。
有向图需要重新明确入度/出度 Laplacian 与根向内/向外树的方向，不能直接复用无向封装。

\section{概率、期望与博弈}\index{expectation}\index{Sprague-Grundy}
期望线性性不要求独立：$\mathbb E[\sum X_i]=\sum\mathbb E[X_i]$。
方差可加需要协方差为零；独立是充分条件。非负整数变量
$\mathbb E[X]=\sum_{k\ge1}\Pr(X\ge k)$。
吸收过程的期望方程 $E_i=1+\sum_jp_{ij}E_j$ 需要检查是否几乎必然吸收及期望是否有限；
不能对闭合非吸收分量直接求唯一有限解。
有限无环公平组合游戏在正常玩法下
\[
\operatorname{SG}(u)=\operatorname{mex}\{\operatorname{SG}(v):u\to v\}.
\]
独立子游戏的 SG 值异或为零当且仅当先手必败。
反常玩法、含环游戏、双方可走步不同的游戏不能直接套这个结论。

\section{数值算法与几何}\index{numerical stability}\index{Pick theorem}
排序必须使用严格弱序；不能以“差小于 eps 则相等”定义通用排序比较器。
叉积与行列式容易发生抵消；整数题优先精确整数判定，浮点构造单独处理误差。
对于简单格点多边形（无孔），Pick 定理 $A=I+B/2-1$，边界格点数
$B=\sum\gcd(|\Delta x|,|\Delta y|)$；不要重复计入边端点。
自适应积分的局部误差估计需要函数足够规则；奇点、间断与强振荡必须先分段处理。
三分搜索要求单峰性，整数域最后枚举剩余区间；二分答案要求判定单调性。

\paragraph{来源索引}
本章为重新组织的公式速查，扩展范围依据 \url{https://oi-wiki.org/math/}。
小模数 exLucas 与有限域拉格朗日插值已有双风格实现和本地验证；素数模平方根已有 Cipolla 实现；仍需补齐二次剩余的其他方法与定理、原根、Min25、Newton 插值、子集卷积、行列式进阶、数值积分等内容及验证。本章公式不能代替这些待完成代码。

\section{线性递推：系数方向与外推条件}
\index{线性递推}
约定 $a_i=\sum_{j=1}^{k}c_{j-1}a_{i-j}$，初值是
$a_0,\ldots,a_{k-1}$。对应特征多项式为
\[
F(x)=x^k-c_0x^{k-1}-c_1x^{k-2}-\cdots-c_{k-1}.
\]
若 $x^n\bmod F(x)=\sum_{i=0}^{k-1}r_ix^i$，则
$a_n=\sum_{i=0}^{k-1}r_ia_i$。代码由高次向低次消去多项式项，
第 $j$ 个系数写入原次数减 $j$ 的位置，避免倒置系数。
例如 Fibonacci 初值为 $[0,1]$，递推系数为 $[1,1]$。

BM 在域上求有限前缀的最短递推。本书实现要求素数模数；
一般的模合数环不能保证非零偏差存在逆元。
仅凭若干样本无法证明真实序列以后一直遵循同一个规律。
若事先知道序列满足阶数至多 $d$ 的齐次线性递推，
至少给出前 $2d$ 项，再用恢复出的系数求后续项。
系数向量末尾的零不能擅自删除：它可能对应非平稳的初始段，
改变阶数会改变递推开始适用的位置。
